\subsection{INVARIANT ELEMENTS}

Let \(A[P]^W\) be the subalgebra of A{[}P{]} consisting of the elements invariant under W. For \(p \in P\) , denote by W.p the orbit of p under W, and let \(S(e^p) = \sum_{q \in W \cdot p} e^q\) be the sum of the W-transforms of \(e^p\) ; this is a W-invariant element. If \(p \in P \cap \overline{C}\) , \(w(p) \leqslant p\) for all \(w \in W\) (§ 1, no. 6, Prop. 18) and \(e^p\) is the unique maximal term of \(S(e^p)\) .

Let \(x = \sum_{p} x_{p} e^{p} \in \mathrm{A}[\mathrm{P}]^{\mathbb{W}}\) ; then \(x_{w(p)} = x_{p}\) for all \(p \in \mathrm{P}\) and all \(w \in \mathrm{W}\) . On the other hand, every orbit of \(\mathbb{W}\) in \(\mathbb{P}\) meets \(\mathbb{P} \cap \overline{\mathbb{C}}\) in exactly one point (§ 1, no. 5, Th. 2). Hence,

\[
x = \sum_ {P \cap \overline {{C}}} x _ {p} S (e ^ {p}).\tag{6}
\]

We deduce:

Lemma 3. The \(S(e^p)\) for \(p \in P \cap \overline{C}\) form a basis of the A-module \(A[P]^W\) .

More generally:

PROPOSITION 3. For any \(p \in \mathrm{P} \cap \overline{\mathrm{C}}\) , let \(x_p\) be an element of \(\mathrm{A}[\mathrm{P}]^{\mathrm{W}}\) with unique maximal term \(e^p\) . Then the family \((x_p)_{p \in \mathrm{P} \cap \overline{\mathrm{C}}}\) is a basis of the A-module \(\mathrm{A}[\mathrm{P}]^{\mathrm{W}}\) .

We first prove a lemma.

Lemma 4. Let I be an ordered set satisfying the following condition:

(MIN) Every non-empty subset of I contains a minimal element.

Let \(\mathbf{E}\) be an A-module, \((e_i)_{i \in I}\) a basis of \(\mathbf{E}\) and \((x_i)_{i \in I}\) a family of elements of \(\mathbf{E}\) such that

\[
x _ {i} = e _ {i} + \sum_ {j <   i} a _ {i j} e _ {j},
\]

for all \(i \in \mathbf{I}\) (with \(a_{ij} \in \mathbf{A}\) , the support of the family \((a_{ij})\) being finite for all \(i\) ). Then, \((x_i)_{i \in \mathbf{I}}\) is a basis of \(\mathbf{E}\) .

For any subset \(J\) of \(I\) , let \(E_J\) be the submodule of \(E\) with basis \((e_i)_{i \in J}\) . Let \(S\) be the set of subsets \(J\) of \(I\) with the following properties:

\begin{enumerate}
\def\labelenumi{(\alph{enumi})}
\item
  If \(i' \leqslant i\) and \(i \in \mathbf{J}\) , then \(i' \in \mathbf{J}\) ;
\item
  \((x_{i})_{i\in J}\) is a basis of \(\mathrm{E_J}\) .
\end{enumerate}

It is immediate that S, ordered by inclusion, is inductive and non-empty. It therefore has a maximal element J. If \(J \neq I\) , let \(i_{0}\) be a minimal element of I - J and put \(J' = J \cup \{i_{0}\}\) . Every element \(i \in I\) such that \(i < i_{0}\) then belongs to J: it follows that \(J'\) satisfies (a). On the other hand, \(J'\) also satisfies (b): indeed,

\[
e _ {i _ {0}} = x _ {i _ {0}} - \sum_ {j <   i _ {0}} a _ {i _ {0} j} e _ {j},
\]

from which (b) follows. Hence \(J' \in S\) , a contradiction. Thus, J = I and the lemma is proved.

We now prove Prop. 3. We apply Lemma 4 with \(\mathbf{I} = \mathbf{P} \cap \overline{\mathbf{C}}\) . Let \(q \in \mathbf{I}\) , and let \(\mathbf{I}_q\) be the set of \(p \in \mathbf{I}\) such that \(p \leqslant q\) . If \(p \in \mathbf{I}_q\) , the relations

\[
q - p \geqslant 0, \quad p \in \overline {{\mathrm{C}}}, \quad q \in \overline {{\mathrm{C}}}
\]

imply that

\[
(q - p | p) \geqslant 0 \text { and } (q - p | q) \geqslant 0,
\]

and hence that

\[
(p | p) \leqslant (p | q) \leqslant (q | q).
\]

The set \(I_{q}\) is thus bounded. Since I is discrete, it follows that \(I_{q}\) is finite, and it is clear that I satisfies the condition (MIN). On the other hand, for all \(p \in I\) ,

\[
x _ {p} = e ^ {p} + \sum_ {q <   p} c _ {p q} e ^ {q}
\]

so by (6).

\[
x _ {p} = \mathrm{S} (e ^ {p}) + \sum_ {q <   p, q \in \mathbb {I}} c _ {p q} \mathrm{S} (e ^ {q}).
\]

The proposition now follows from Lemmas 3 and 4.

THEOREM 1. Let \(\overline{\omega}_1, \ldots, \overline{\omega}_l\) be the fundamental weights corresponding to the chamber C, and, for \(1 \leqslant i \leqslant l\) , let \(x_i\) be an element of A{[}P{]} \(^W\) with \(e^{\overline{\omega}_i}\) as its unique maximal term. Let

\[
\varphi : \Lambda [ X _ {1}, \dots , X _ {I} ] \rightarrow A [ P ] ^ {W}
\]

be the homomorphism from the polynomial algebra \(A[X_{1},\ldots,X_{l}]\) to \(A[P]^{W}\) that takes \(X_{i}\) to \(x_{i}\) . Then, the map \(\varphi\) is an isomorphism.

Lemma 2 implies that the image under \(\varphi\) of the monomial \(X_{1}^{n_{1}}\ldots X_{l}^{n_{l}}\) is an element with unique maximal term \(e^{n_1\overline{\omega}_1 + \dots +n_l\overline{\omega}_l}\) . Since every element of \(P\cap \overline{C}\) can be written uniquely in the form \(n_1\overline{\omega}_1 + \dots +n_l\overline{\omega}_l\) , Prop. 3 shows that the images under \(\varphi\) of the monomials \(X_{1}^{n_{1}}\ldots X_{l}^{n_{l}}\) are a basis of \(A[P]^W\) , hence the theorem.

Examples. 1) We can take \(x_{i} = \mathrm{S}(\mathfrak{e}^{\overline{\omega}_i})\) .

\begin{enumerate}
\def\labelenumi{\arabic{enumi})}
\setcounter{enumi}{1}
\tightlist
\item
  By Remark 2 of no. 3. we can take \(x_{i} = \mathrm{J}(e^{\rho +\overline{\omega}_i}) / d\) (with the notation in no. 3).
\end{enumerate}

